% GENERATED FILE. Edit canonical lessons, then run npm run generate:books.
% canonical-source-hash: fnv1a64:780be4eb7eee8087
% canonical-lessons: GU-C174-oktobar, GU-C174-disembar, GU-C174-maha, GU-C174-tarikhiyu, GU-C174-saruat

\chapter{October to Maha}
\label{ch:gu-october-to-maha}
\hlchaptermodality{174}

\begin{chapteropening}
\textbf{By the end of this chapter:} \emph{I can say in which month something will happen.}

Complete the last lesson of chapter 174: I can say in which month something will happen.
\end{chapteropening}

\section[\texorpdfstring{\gu{ઓક્ટોબર} (okṭobar)}{ઓક્ટોબર (okṭobar)}]{\gu{ઓક્ટોબર} (okṭobar) — October}

\label{lesson:GU-C174-oktobar}



\begin{quote}
Before the new one: say the Gujarati for last, then the Gujarati for whole.
\end{quote}

\subsection*{You'll want to know: \gu{ઓક્ટોબર}}
\textbf{\gu{ઓક્ટોબર}} — \emph{okṭobar} — \textquotedblleft{}October\textquotedblright{}.

\begin{quote}
\textbf{\gu{ઓક્ટોબરમાં} \gu{દિવાળી} \gu{આવશે}.} — \emph{okṭobarmā̃ divāḷī āvśe} — \textquotedblleft{}Diwali will come in October.\textquotedblright{}
\end{quote}

\begin{cousinweb}
\textbf{\gu{ઓક્ટોબર}} (\emph{okṭobar}) is a loan from English \emph{October}, the name of the tenth month.
\end{cousinweb}

\begin{grammarlens}[title={What you've built}]
One more word for everyday life at A2.
\end{grammarlens}

\subsection*{Your turn}
\begin{itemize}
\raggedright
  \item \emph{Say it:} \emph{okṭobar}
  \item \emph{Say it:} \emph{okṭobar}, once more
  \item \emph{Say it:} \emph{okṭobarmā̃ divāḷī āvśe} — Diwali will come in October.
  \item \emph{Recall:} say the Gujarati for last, then the Gujarati for whole, then say \emph{okṭobar} again
\end{itemize}

\begin{tcolorbox}[breakable,colback=teal!4,colframe=teal!35!black,title={Before you move on}]
What does \gu{ઓક્ટોબર} mean? (\textquotedblleft{}October\textquotedblright{}.) Say it once more.
\end{tcolorbox}
\section[\texorpdfstring{\gu{ડિસેમ્બર} (ḍisembar)}{ડિસેમ્બર (ḍisembar)}]{\gu{ડિસેમ્બર} (ḍisembar) — December}

\label{lesson:GU-C174-disembar}



\begin{quote}
Before the new one: say the Gujarati for whole, then the Gujarati for October.
\end{quote}

\subsection*{You'll want to know: \gu{ડિસેમ્બર}}
\textbf{\gu{ડિસેમ્બર}} — \emph{ḍisembar} — \textquotedblleft{}December\textquotedblright{}.

\begin{quote}
\textbf{\gu{ડિસેમ્બરમાં} \gu{શિયાળો} \gu{હશે}.} — \emph{ḍisembarmā̃ śiyāḷo haśe} — \textquotedblleft{}It will be winter in December.\textquotedblright{}
\end{quote}

\begin{cousinweb}
\textbf{\gu{ડિસેમ્બર}} (\emph{ḍisembar}) is a loan from English \emph{December}, the name of the twelfth month.
\end{cousinweb}

\begin{grammarlens}[title={What you've built}]
One more word for everyday life at A2.
\end{grammarlens}

\subsection*{Your turn}
\begin{itemize}
\raggedright
  \item \emph{Say it:} \emph{ḍisembar}
  \item \emph{Say it:} \emph{ḍisembar}, once more
  \item \emph{Say it:} \emph{ḍisembarmā̃ śiyāḷo haśe} — It will be winter in December.
  \item \emph{Recall:} say the Gujarati for whole, then the Gujarati for October, then say \emph{ḍisembar} again
\end{itemize}

\begin{tcolorbox}[breakable,colback=teal!4,colframe=teal!35!black,title={Before you move on}]
What does \gu{ડિસેમ્બર} mean? (\textquotedblleft{}December\textquotedblright{}.) Say it once more.
\end{tcolorbox}
\section[\texorpdfstring{\gu{મહા} (mahā)}{મહા (mahā)}]{\gu{મહા} (mahā) — Maha, a winter month}

\label{lesson:GU-C174-maha}



\begin{quote}
Before the new one: say the Gujarati for October, then the Gujarati for December.
\end{quote}

\subsection*{You'll want to know: \gu{મહા}}
\textbf{\gu{મહા}} — \emph{mahā} — \textquotedblleft{}Maha, a winter month\textquotedblright{}.

\begin{quote}
\textbf{\gu{મહા} \gu{મહિનામાં} \gu{લગ્ન} \gu{હશે}.} — \emph{mahā mahināmā̃ lagna haśe} — \textquotedblleft{}The wedding will be in the month of Maha.\textquotedblright{}
\end{quote}

\begin{cousinweb}
\textbf{\gu{મહા}} (\emph{mahā}) is a tadbhava. It comes from Sanskrit \emph{māgha}, \textquotedblleft{}the month January-February\textquotedblright{}, and CDIAL lists the Gujarati \emph{mahā}, \emph{māh} for that month (CDIAL 9993).
\end{cousinweb}

\begin{grammarlens}[title={What you've built}]
One more word for everyday life at A2.
\end{grammarlens}

\subsection*{Your turn}
\begin{itemize}
\raggedright
  \item \emph{Say it:} \emph{mahā}
  \item \emph{Say it:} \emph{mahā}, once more
  \item \emph{Say it:} \emph{mahā mahināmā̃ lagna haśe} — The wedding will be in the month of Maha.
  \item \emph{Recall:} say the Gujarati for October, then the Gujarati for December, then say \emph{mahā} again
\end{itemize}

\begin{tcolorbox}[breakable,colback=teal!4,colframe=teal!35!black,title={Before you move on}]
What does \gu{મહા} mean? (\textquotedblleft{}Maha, a winter month\textquotedblright{}.) Say it once more.
\end{tcolorbox}
\section[\texorpdfstring{\gu{તારીખિયું} (tārīkhiyũ)}{તારીખિયું (tārīkhiyũ)}]{\gu{તારીખિયું} (tārīkhiyũ) — wall calendar}

\label{lesson:GU-C174-tarikhiyu}



\begin{quote}
Before the new one: say the Gujarati for December, then the Gujarati for Maha.
\end{quote}

\subsection*{You'll want to know: \gu{તારીખિયું}}
\textbf{\gu{તારીખિયું}} — \emph{tārīkhiyũ} — \textquotedblleft{}wall calendar\textquotedblright{}.

\begin{quote}
\textbf{\gu{હું} \gu{નવું} \gu{તારીખિયું} \gu{ખરીદીશ}.} — \emph{hũ navũ tārīkhiyũ kharīdīś} — \textquotedblleft{}I will buy a new calendar.\textquotedblright{}
\end{quote}

Built on \textbf{\gu{તારીખ}}, \emph{tārīkh}, date.

\begin{cousinweb}
\textbf{\gu{તારીખિયું}} (\emph{tārīkhiyũ}) is built on the taught \emph{tārīkh}, \textquotedblleft{}date\textquotedblright{}, with the ending \emph{-iyũ}. \emph{Tārīkh} is a loan from Arabic \emph{tārīkh}, \textquotedblleft{}a dating, a date\textquotedblright{}, through Persian; EtymDB lists the Persian word. Its Arabic root is \emph{ʾ-r-kh}, \textquotedblleft{}to date a letter\textquotedblright{} (Lane).
\end{cousinweb}

\begin{grammarlens}[title={What you've built}]
One more word for everyday life at A2.
\end{grammarlens}

\subsection*{Your turn}
\begin{itemize}
\raggedright
  \item \emph{Say it:} \emph{tārīkhiyũ}
  \item \emph{Say it:} \emph{tārīkhiyũ}, once more
  \item \emph{Say it:} \emph{hũ navũ tārīkhiyũ kharīdīś} — I will buy a new calendar.
  \item \emph{Recall:} say the Gujarati for December, then the Gujarati for Maha, then say \emph{tārīkhiyũ} again
\end{itemize}

\begin{tcolorbox}[breakable,colback=teal!4,colframe=teal!35!black,title={Before you move on}]
What does \gu{તારીખિયું} mean? (\textquotedblleft{}Wall calendar\textquotedblright{}.) Say it once more.
\end{tcolorbox}
\section[\texorpdfstring{\gu{શરૂઆત} (śarūāt)}{શરૂઆત (śarūāt)}]{\gu{શરૂઆત} (śarūāt) — beginning, start}

\label{lesson:GU-C174-saruat}



\begin{quote}
Before the new one: say the Gujarati for Maha, then the Gujarati for wall calendar.
\end{quote}

\subsection*{You'll want to know: \gu{શરૂઆત}}
\textbf{\gu{શરૂઆત}} — \emph{śarūāt} — \textquotedblleft{}beginning, start\textquotedblright{}.

\begin{quote}
\textbf{\gu{મહિનાની} \gu{શરૂઆતમાં} \gu{પગાર} \gu{આવશે}.} — \emph{mahinānī śarūātmā̃ pagār āvśe} — \textquotedblleft{}The salary will come at the start of the month.\textquotedblright{}
\end{quote}

\begin{cousinweb}
\textbf{\gu{શરૂઆત}} (\emph{śarūāt}) is a loan from Arabic \emph{shurūʿ}, \textquotedblleft{}a beginning\textquotedblright{}, which reached Gujarati through Persian, with the ending \emph{-āt}; EtymDB lists the Arabic word. Its root is \emph{sh-r-ʿ}, \textquotedblleft{}to enter upon\textquotedblright{} (Lane).
\end{cousinweb}

\begin{grammarlens}[title={What you've built}]
One more word for everyday life at A2.
\end{grammarlens}

\subsection*{Your turn}
\begin{itemize}
\raggedright
  \item \emph{Say it:} \emph{śarūāt}
  \item \emph{Say it:} \emph{śarūāt}, once more
  \item \emph{Say it:} \emph{mahinānī śarūātmā̃ pagār āvśe} — The salary will come at the start of the month.
  \item \emph{Recall:} say the Gujarati for Maha, then the Gujarati for wall calendar, then say \emph{śarūāt} again
\end{itemize}

\begin{tcolorbox}[breakable,colback=teal!4,colframe=teal!35!black,title={Before you move on}]
What does \gu{શરૂઆત} mean? (\textquotedblleft{}Beginning, start\textquotedblright{}.) Say it once more.
\end{tcolorbox}
